\documentclass[12pt,a4paper]{article}
\usepackage{minekommandoer}
\begin{document}
\begin{center}
    \setfontsize{16pt}{\textbf{Prøve i støkiometri:}}
\end{center}\vspace{1cm}
\textbf{Oppg. 1:}\nl
a) beregn den molare massen til følgende stoffer:\nl
$H_2O_2$, $NH_3$, $H_2SO_4$,$CaCl_2$,$KMnO_4$,$Ca(NO_3)_2$\nl
b) Beregn hvor mange mol du har av stoffene i a når du har 24,5g av stoffene.\nl
c) Beregn hvor mange gram du har av stoffene i a når du har 34,5mol av stoffene.\nl
\textbf{Oppg. 2:}\nl
Gitt følgende kjemiske reaksjon:\nl
$$C+O_2=CO_2$$\nl
a) Finn molar mase til C, $O_2$ og $CO_2$.\nl
b) Hvor mye får vi laget av $CO_2$ og hvor mye $O_2$ trenger vi når vi har 36g C?\nl
c) Hvor mange gram trenger vi av C og $O_2$ for å lage 245g $CO_2$?\nl
\newpage
\textbf{Oppg. 3:}\nl
Gitt følgende reaksjon:
$$N_2+3H_2=2NH_3$$
a) hvor mye ammoniakk ($NH_3$) får vi dannet om vi starter med 42g $N_2$, og hvor mye $H_2$ trenger vi da?\nl
b) Hvor mye ammobniakk får du dannet om du startrer med 42g $N_2$ og 5g $H_2$?
\nl[5]
\textbf{Formel for å beregne antall mol:}\nl[2]
$$n=\frac{m}{Mm}$$
Som trekant:\nl
\trekant{m}{n}{Mm}
n=stoffmengde (mol)\nl
m=masse(g)\nl
Mm=molar masse (g/mol)\nl
\end{document} 